%!TEX root = ../../main.tex
\section{물질적 결과와 다음 오브젝트와의 대응}
\label{sec:value-correspondence}

$\SVI$가 같은 사례의 물질적 결과와 어떻게 대응하는지를 결과 시점(h90) 상태에서 계산한 다섯 축(킬 차이, 엘리트 몬스터 순차, 구조물 순차, 둘을 합한 오브젝트 순차, 종점 생존 차이)에 대해 본다(표 \ref{tab:material-axes}; 보완 실험, 15.16 노출 이후의 진단). 킬 차이가 결정된 교전에서 킬 축 부호와 $\SVI$의 경기 가중 일치율은 T 전체에서 \maTKillAllAgree{}, $\Bforty$에서 \maTKillBfortyAgree{}이며 S에서는 \maSKillAllAgree{}과 \maSKillBfortyAgree{}이다. 나머지 축은 동률 비율이 높고(예: T의 엘리트 몬스터 순차 \maTEpicAllTie{}) 일치율은 킬 축보다 낮다. 물질 특징은 $\Vhat$의 입력과 겹치므로 이 표는 같은 사례 안의 대응이다.

%% GENERATED by thesis/tools/gen_supp.py -- do not edit
\begin{table}[!t]
  \centering
  \caption{SVI와 물질적 결과 축의 대응 (15.16 TEST)}
  \label{tab:material-axes}
  \footnotesize

  \begin{tabular}{@{}lllrrrr@{}}
    \toprule
    코호트 & 축 & 범위 & 결정 $n$ & 동률 비율 & 일치율 & 불일치율 \\
    \midrule
    T & 킬 차이 & 전체 & 28{,}711 & 0.129 & 0.904 & 0.097 \\
    T & 킬 차이 & $B_{40}$ & 4{,}653 & 0.142 & 0.928 & 0.074 \\
    T & 엘리트 몬스터 순차 & 전체 & 12{,}507 & 0.621 & 0.818 & 0.182 \\
    T & 엘리트 몬스터 순차 & $B_{40}$ & 2{,}221 & 0.590 & 0.842 & 0.158 \\
    T & 구조물 순차 & 전체 & 16{,}106 & 0.512 & 0.774 & 0.228 \\
    T & 구조물 순차 & $B_{40}$ & 2{,}014 & 0.629 & 0.812 & 0.190 \\
    T & 오브젝트 순차 (엘리트+구조물) & 전체 & 21{,}779 & 0.340 & 0.810 & 0.191 \\
    T & 오브젝트 순차 (엘리트+구조물) & $B_{40}$ & 3{,}099 & 0.429 & 0.860 & 0.141 \\
    T & 종점 생존 차이 & 전체 & 14{,}584 & 0.558 & 0.794 & 0.208 \\
    T & 종점 생존 차이 & $B_{40}$ & 2{,}142 & 0.605 & 0.788 & 0.212 \\
    S & 킬 차이 & 전체 & 83{,}667 & 0.173 & 0.894 & 0.108 \\
    S & 킬 차이 & $B_{40}$ & 25{,}574 & 0.193 & 0.923 & 0.080 \\
    S & 엘리트 몬스터 순차 & 전체 & 23{,}599 & 0.767 & 0.715 & 0.284 \\
    S & 엘리트 몬스터 순차 & $B_{40}$ & 4{,}761 & 0.850 & 0.714 & 0.285 \\
    S & 구조물 순차 & 전체 & 20{,}502 & 0.797 & 0.684 & 0.317 \\
    S & 구조물 순차 & $B_{40}$ & 2{,}508 & 0.921 & 0.719 & 0.281 \\
    S & 오브젝트 순차 (엘리트+구조물) & 전체 & 35{,}907 & 0.645 & 0.709 & 0.292 \\
    S & 오브젝트 순차 (엘리트+구조물) & $B_{40}$ & 6{,}060 & 0.809 & 0.726 & 0.273 \\
    S & 종점 생존 차이 & 전체 & 34{,}992 & 0.654 & 0.759 & 0.243 \\
    S & 종점 생존 차이 & $B_{40}$ & 8{,}345 & 0.737 & 0.782 & 0.219 \\
    \bottomrule
  \end{tabular}
  \tabnote{축 = 결과 시점(h90) 상태에서 계산한 블루$-$레드 순차(킬·엘리트 몬스터·구조물·둘의 합·생존 인원). 결정 = 축이 0이 아닌 행; 일치율 = 결정 행 가운데 SVI 부호와 축 부호가 같은 비율(경기 가중); 불일치율은 비가중. 물질 특징은 $\hat V$의 입력과 겹치므로 같은 사례 안의 대응이다.}
\end{table}


킬 축과 $\SVI$가 다른 사례에서 오브젝트 순차가 어느 쪽과 일치하는지도 세었다. T에서 킬 결정 행 가운데 두 부호가 다른 \kdTN{}건 중 오브젝트 순차는 $\SVI$ 쪽과 \kdTSvi{}건, 킬 쪽과 \kdTKill{}건 일치하고 \kdTTie{}건은 동률이다. S에서는 \kdSN{}건 중 \kdSSvi{} / \kdSKill{} / \kdSTie{}건이다. 두 라벨이 갈리는 사례는 물질 순차로도 한쪽으로 기울지 않으며, 이 분포는 두 결과 정의가 서로 다른 사례를 다르게 평가한다는 사실의 기술이다.

결과 시점 이후 180초 창에서 첫 엘리트 오브젝트(바론, 용, 전령, 유충, 아타칸)가 한 팀에 귀속된 경우 가운데, 블루에 귀속된 비율은 $\SVI+$ 뒤에서 \objBluePos{}(\objBluePosNum{} / \objBluePosDen{}), $\SVI-$ 뒤에서 \objBlueNeg{}(\objBlueNegNum{} / \objBlueNegDen{})이다(표 \ref{tab:correspondence}; 비가중 건수). 창 안에서 경기가 끝나 오브젝트가 없는 경우가 \objGameEnded{} / \numTTest{}이고 관측 중단은 0이다. 같은 라벨을 다른 분모로 읽으면(표 \ref{tab:nextobj}) 귀속 팀이 Blue 또는 Red인 \noAllNBR{}행에서 $\SVI$ 방향과 귀속 팀이 일치한 비율 $P(\Ysvi = \ind[O = \mathrm{Blue}] \mid O \in \{\mathrm{Blue}, \mathrm{Red}\})$은 \noAllSviAgree{}이다. 이 값은 위의 조건부 귀속 비율 \objBluePos{}와 분자·분모가 다른 통계량이며 반올림한 값이 같을 뿐이다. 같은 행 가운데 킬 차이가 0이 아닌 \noAllNKill{}행에서 킬 부호와 귀속 팀의 일치율은 \noAllKillAgree{}이고, 두 일치율은 분모가 다르므로 같은 척도 위에 있지 않다.

%% table 객체: tab:correspondence — 출처 RR5_RR6B (RR5a crosstab; RR5b L35–38)
\begin{table}[htbp]
  \centering
  \caption{$\SVI$와 관측 결과의 대응 (T, 15.16). RR5a는 경기 가중 일치율, RR5b는 비가중 건수. 대응은 $q$의 정확도가 아니다.}
  \label{tab:correspondence}
  \small
  \begin{tabular}{@{}llr@{}}
    \toprule
    검사 & 셀 & 값 \\
    \midrule
    킬 축 부호 일치 (결정된 경우) & 전체 T & \killAgreeAll{} \\
    킬 축 부호 일치 (결정된 경우) & $\Bforty$ & \killAgreeBforty{} \\
    180 s 내 첫 엘리트 오브젝트의 블루 귀속 비율 & $\SVI+$ 뒤 & \objBluePos{} (\objBluePosNum{} / \objBluePosDen{}) \\
    180 s 내 첫 엘리트 오브젝트의 블루 귀속 비율 & $\SVI-$ 뒤 & \objBlueNeg{} (\objBlueNegNum{} / \objBlueNegDen{}) \\
    창 안 경기 종료 (오브젝트 없음) & 전체 T & \objGameEnded{} / \numTTest{} \\
    \bottomrule
  \end{tabular}
\end{table}

%% GENERATED by thesis/tools/gen_supp.py -- do not edit
\begin{table}[!t]
  \centering
  \caption{다음 엘리트 오브젝트와의 연관 (한타 T, 15.16 TEST)}
  \label{tab:nextobj}
  \footnotesize

  \begin{tabular}{@{}lrrrrrrr@{}}
    \toprule
    범위 & $n$ & Blue & Red & Blue$|$Red & SVI 일치율 & 킬 결정 & 킬 일치율 \\
    \midrule
    전체 & 32{,}981 & 10{,}076 & 10{,}135 & 20{,}211 & 0.604 & 17{,}053 & 0.670 \\
    $B_{40}$ & 5{,}423 & 1{,}894 & 1{,}869 & 3{,}763 & 0.610 & 3{,}158 & 0.638 \\
    SVI $=1$ & 16{,}358 & 6{,}004 & 3{,}930 & 9{,}934 & 0.604 & 8{,}378 & 0.664 \\
    SVI $=0$ & 16{,}623 & 4{,}072 & 6{,}205 & 10{,}277 & 0.604 & 8{,}675 & 0.676 \\
    \bottomrule
  \end{tabular}
  \tabnote{$O$ = 결과 시점 뒤 180 s 안의 첫 엘리트 오브젝트 획득 팀(RR5b 라벨 재사용). Blue$|$Red = $O\in\{\mathrm{Blue},\mathrm{Red}\}$인 행(없음·경기 종료·동률 제외). SVI 일치율 = $P(Y_{\mathrm{SVI}} = \mathbf 1[O=\mathrm{Blue}] \mid \mathrm{Blue}|\mathrm{Red})$; 킬 일치율 = 같은 조건에 킬 차이 $\ne 0$을 더한 분모에서 킬 부호와 $O$의 일치율. 두 일치율은 분모가 다르다. 비가중, 구간 없음.}
\end{table}

